\section{Introduction}
A reusable description specifies one quantum instrument whose behaviour
remains close to the supplied target through every allowed interaction.
This problem depends not only on how many real parameters the target
has, but on how repeated use resolves its different directions.
We study that distinction for varying projective readout and singular
conditional preparations.

For two memoryless instruments of the same interface, define
\[
 d_N(\Phi,\Psi)=\sup_{\mathsf T}
       \|\mathsf T[\Phi]-\mathsf T[\Psi]\|_1.
\]
The same causal tester $\mathsf T$ is used in both hypotheses. It may
have initial reference entanglement, arbitrary finite quantum memory,
feedback and a public stopping time bounded by $N$. The output includes
its retained classical and quantum systems. Thus the distance is
unhalved and has maximum two; event probabilities differ by at most
$d_N/2$. Each device invocation uses the same memoryless instrument.
Covering centres are legal instruments, not arbitrary arrays.

Our first family is
\[
 \mathcal F_{P,\sigma,(x,y)}(X)=\tr(P_xX)\sigma_{xy}.
\]
Here the $P_x$ form an ordered rank-one projective measurement, both
labels $(x,y)$ are retained, and each conditional row of positive
blocks is trace normalized. The rank bounds $r_{xy}$ allow changing
supports and zero blocks. Put $b=d^2-d$ and
$V=\sum_x(\sum_y r_{xy}(2n-r_{xy})-1)$.
Theorem~\ref{thm:readoutentropy72} proves
\[
 \log_2\mathcal R_N(\delta)
  =(b+V/2)\log_2N+(b+V)\log_2(1/\delta)+O(1),
 \qquad 0<\delta\le1/32.
\]
The remainder is independent of $N\ge1$ and $\delta$, at fixed
public dimensions and ranks. The lower bound allows arbitrary legal
memoryless centres; the upper code returns rational family members
without increasing conditional ranks. This is an observable-readout
theorem, not a dimension count for all entanglement-breaking channels.
Corollary~\ref{cor:observablereadout72} removes the explicit $x$ flag
when fixed orthogonal output supports recover it. Equal overlapping
unlabelled rows, by contrast, can erase every basis parameter.

The basis scale follows from the established multiple-query theorem
for von Neumann measurements \cite{PPKKO72}, not from the intrinsic
dimension of the flag manifold alone. Readout differences remain
visible when the conditional outputs are discarded. Conditional-state
differences at a common readout are visible by a basis-state input.
These two pair-dependent witnesses justify a product packing at scales
$\delta/N$ and $\delta/\sqrt N$. For the upper code, partial-pivot
elimination selects bounded rational coordinates directly from the
ordered projections. Normalized algebraic eigenvectors and exhaustive
unitary-net search are unnecessary. Positive row-factor codes and a
common programme comparison then control all adaptive testers.

A second result, Theorem~\ref{thm:seizablecover72}, identifies the
covering problem exactly when one common tester extracts the programme
from one use throughout its ambient state domain. State/channel
comparisons under environment parametrization and seizing are classical
\cite{DW72,WBHK72,WW72}. The additional centre argument retracts every
legal memoryless centre through the same seizer and processor. Thus
finite-use family covering numbers equal their state counterparts,
not merely a binary asymptotic exponent. The inherited rank-state law
then gives the full range $0<\delta\le\delta_*<2$ for fixed cap;
Corollary~\ref{cor:paulicoding72} illustrates it with Pauli channels.
The retraction may increase Choi rank. Rank and zero preservation
belong to the upper target encoder, not to arbitrary-centre retraction.

The preceding results are placed alongside their necessary preliminary
geometry. The fixed-readout input-dependent theorem
\ref{thm:cqentropy71} has coefficient $V/2$ in $\log N$ and the full
fixed submaximal-error range. The common-estimator theorem
\ref{thm:fullpacking71} explains that range by disjoint-event
multiplicity, rather than pairwise separation beyond half the diameter.
The coherent/preparation tensor theorem
\ref{thm:coherententropy70} has projective-unitary dimension $u=d^2-1$
and coefficient $u+v/2$; its small-error range is unchanged.
These earlier theorems, the intrinsic Choi codec and the singular
preparation code are retained with complete proofs. In particular the
new readout quotient removes all column phases, not merely one global
phase. Revision genealogy is recorded outside the article.

The encoder receives matrix data. Fixed dimensions, rank bounds,
horizon, tolerance and error cap are public. Finite chart headers,
pivot masks and coordinate digits are charged in the reusable payload;
expanded matrices, encoder workspace and quantum programme hardware
are different resources. Exact rational implementation applies to
rational targets; the real-target theorem constructs rational centres
in exact real arithmetic. No theorem here learns an unknown device
from queries or implements its action with a finite classical message.
